\chapter{Introduction}
\section{Background and problem}
Crop pathogens and pests impose substantial losses, with effects that vary by crop and geography~\cite{savary2019losses}. Official Indian horticulture statistics establish tomato as a substantial production crop~\cite{pib2026tomato}. Leaf imagery therefore offers a useful application for accessible visual screening. PlantVillage studies demonstrate classification on curated images~\cite{hughes2015dataset,mohanty2016plant}; they also motivate careful separation of benchmark performance from field reliability.

This project asks whether Differential Evolution (DE) can improve an ImageNet-pretrained EfficientNet-B0 tomato classifier by selecting its training recipe. The objective is validation macro-F1 over a mixed space of learning rate, regularization, head width, fine-tuning depth, and augmentation. Backpropagation fits each candidate's weights; DE selects the recipe. The dataset, backbone family, and optimizer topic remain fixed.

\section{Measurable objectives}
\begin{enumerate}
\item Prepare documented class mappings, cleaning records, and group-aware stratified splits.
\item Implement and test DE donor selection, mutation, repair, crossover, and selection.
\item Compare a registered baseline, DE, and equal-budget random search using validation-only selection.
\item Evaluate repeated seeds, per-class performance, uncertainty, robustness, and calibration.
\item Deliver an offline demonstration and reproducible report, slides, code, and evidence.
\end{enumerate}

\section{Contribution, scope, and organization}
The contribution is an auditable chain from candidate vector to checkpoint and reported result. Established methods are combined and evaluated; no claim of a new optimizer or universal superiority is made. Curated-dataset accuracy does not establish reliability for arbitrary crops, unknown diseases, or field photographs.

Chapters 2--4 develop the literature and algorithms; Chapters 5--6 specify implementation and evaluation. Chapters 7--8 analyze the evidence. Applications, conclusions, future work, verified references, and repository instructions complete the report.
